% ---------------------------------------------------------------------
%  Conclusion générale
% ---------------------------------------------------------------------
\pagetitre{CONCLUSION GÉNÉRALE}

Ce projet nous a permis de concevoir et de réaliser \equipay, une application mobile de
partage des dépenses de groupe qui fait de l'intelligence artificielle un outil concret
plutôt qu'une vitrine. Partant d'un constat simple (la saisie manuelle décourage les
utilisateurs, et rien ne relie les dépenses d'un voyage à son organisation), nous avons
construit une application complète : Flutter pour le mobile, Supabase pour les données
partagées et sécurisées, un backend FastAPI pour les modèles, et MLflow pour suivre leur
cycle de vie.

Sur le plan de l'IA, quatre fonctions ont été intégrées. L'assistant conversationnel, affiné
à partir de Qwen3-4B sur un jeu de données que nous avons construit, passe d'un score de
0,276 à 0,770 et a appris à refuser les questions hors sujet. Le scanner de tickets, affiné à
partir de Qwen3-VL-4B, trouve le bon total sur 97,2 \% des tickets de test et sur la totalité
des tickets en dirhams et en euros. La saisie par phrase atteint un score de 0,976. Enfin,
l'organisation du trajet combine un calcul d'itinéraire, un algorithme d'ordre de visite et
un modèle de langage chargé uniquement de la rédaction. La conversion en 4 bits avec MLX a
réduit de moitié la mémoire nécessaire, ce qui permet de servir tous ces modèles sur un
simple ordinateur portable.

Au-delà des résultats chiffrés, ce projet nous a appris plusieurs leçons. La première
version du scanner, excellente sur son jeu de test mais inutilisable sur nos propres
tickets, nous a montré qu'un bon score ne vaut que si le jeu de test ressemble à l'usage
réel. La séparation entre ce que calcule un algorithme et ce que rédige un modèle s'est
révélée essentielle pour obtenir des résultats fiables. Et la démarche MLOps (données
versionnées, évaluations comparables, choix d'un champion) a transformé une suite d'essais en
une progression mesurable. Le travail en binôme, avec relecture croisée et décisions prises
à deux, a été un apprentissage directement transposable en entreprise.

L'application reste un prototype validé à petite échelle. Les perspectives sont claires :
héberger le backend sur un serveur GPU pour rendre l'IA disponible partout, publier
l'application sur les magasins d'applications (ce qui imposera la connexion avec un compte
Apple), compléter la traduction en anglais et en arabe déjà prévue dans les préférences,
enrichir le jeu de données de l'assistant pour réduire ses réponses inventées sur les
questions de support, et ajouter les notifications. Au terme de ce projet, l'objectif initial
est atteint : \equipay montre qu'il est possible d'intégrer des modèles de langage affinés
dans une application du quotidien, de façon mesurée, traçable et utile.
